% =============================================================================
%  MDIE: Modification-Disentangled Identity Encoder for Occlusion- and
%  Lighting-Robust Face Recognition.
%  Paper draft (IEEEtran).  Numbers reflect the run reported in
%  results/stage2_metrics.json, results/security_family_summary.json,
%  results/real_benchmarks.{csv,json}, results/inference_compat_proof.json,
%  and figures/attention_bone_iou.png  (LFW, min_faces_per_person=8,
%  173 train / 44 held-out identities, identity-disjoint 80/20 split,
%  6 000 verification pairs / modification; single 512-d fused embedding).
% =============================================================================
\documentclass[10pt,conference]{IEEEtran}
\usepackage{graphicx,amsmath,amssymb,booktabs,multirow,url,hyperref,xcolor}
\graphicspath{{../figures/}}

\title{MDIE: Modification-Disentangled Identity Encoding for\\
       Occlusion- and Lighting-Robust Face Recognition}
\author{\IEEEauthorblockN{Anonymous Submission}\\
        \IEEEauthorblockA{Manuscript under review}}

\begin{document}
\maketitle

% =============================================================================
\begin{abstract}
Face encoders trained with margin-based softmax losses (ArcFace, CosFace,
FaceNet, MobileFaceNet) are tuned for cooperative, well-lit capture, yet
\emph{security and access-control} deployments must recognise subjects who are
partially disguised --- surgical masks, caps/hats, eyeglasses, partial
occluders --- under adverse lighting. We show that, in a comparably-trained
regime, these encoders collapse to $0.60$--$0.75$ pooled verification AUC on
such conditions. We propose \textbf{MDIE} (Modification-Disentangled Identity
Encoder), which augments a standard IR-50/ArcFace backbone with three training
signals: (i) \textbf{Adversarial Modification Disentanglement} (AMD), a
gradient-reversal head that removes occlusion/lighting-class information from
the embedding; (ii) an \textbf{Inter-Condition Consistency Loss} (ICCL) with
condition-aware hard mining that pulls clean--modified pairs of the same
identity together; and (iii) \textbf{Region-Aware Token Attention} (RATA), a
spatial attention map supervised onto each face's own detected rigid-bone
landmarks. At inference a learned fusion head merges the bone-anchored
attention embedding and the native identity embedding into a \emph{single
$512$-d, L2-normalised vector produced in one forward pass} --- a drop-in
replacement for an ArcFace encoder under plain cosine matching, with no
modification label required. On a held-out (unseen-identity) LFW protocol with
nine occlusion/lighting modifications, MDIE attains \textbf{pooled AUC $0.979$}
(occlusion family $0.975$, lighting family $0.980$) versus $0.749$ for the best
comparably-trained baseline, and transfers to \emph{real} worn occlusions
(MFR2 masks AUC $0.826$; MeGlass glasses AUC $0.926$), winning every
comparably-trained comparison. We further introduce an \textbf{attention--bone
IoU} interpretability metric with controls: the learned attention overlaps each
individual's \emph{own} bones (IoU $0.694$) far more than a different face's
bones ($0.280$) or a random map ($0.077$), $p\!=\!1.6\!\times\!10^{-15}$ ---
direct evidence that recognition rides on the rigid skeletal scaffold that
survives the disguise. The system trains end-to-end on a single NVIDIA
RTX\,3050 (4\,GB) and is fully reproducible.
\end{abstract}

% =============================================================================
\section{Introduction}
\label{sec:intro}

Margin-based softmax losses
\cite{deng2019arcface,wang2018cosface,schroff2015facenet,liu2017sphereface,zhang2019adacos}
push face-verification benchmarks close to saturation under benign conditions.
Security and surveillance systems, however, operate at the opposite extreme:
subjects are \emph{uncooperative}, frequently wear \emph{disguises} (surgical
masks, caps, sunglasses) or are \emph{partially occluded}, and are lit by
whatever ambient illumination is available (low light, blown-out highlights,
harsh directional shadow). These are exactly the nuisance factors that
margin-based identity objectives never penalise: any cue correlated with
identity in the cooperative training distribution is preserved, so when the
nuisance shifts at test time the embedding rotates along nuisance-aligned
directions and cosine similarity collapses.

We target this operating point directly. We deliberately scope the problem to
\emph{worn occlusions and adverse lighting} for security recognition --- a
sharp, demonstrable niche --- rather than open-ended ``appearance change.''
This paper makes four contributions:
\begin{enumerate}
\item A reproducible occlusion/lighting failure-mode benchmark on top of LFW
      with nine protocol-faithful generators (eyeglasses, surgical mask,
      cap, random occluder, low-light, over-exposure, harsh shadow, aging,
      and FGSM-adversarial, plus a clean control), with identity-disjoint
      pair sampling.
\item \textbf{MDIE}, a training recipe that combines ArcFace identity
      supervision with AMD (adversarial disentanglement), ICCL
      (consistency contrastive loss) and RATA (bone-anchored attention),
      and a fusion head that yields a \emph{single $512$-d
      ArcFace-compatible embedding} at inference.
\item A held-out evaluation showing MDIE lifts pooled occlusion/lighting AUC
      from the $0.749$ baseline regime to $0.979$, and transfers to
      \emph{real} masked (MFR2) and eyeglass (MeGlass) benchmarks, beating
      every comparably-trained baseline.
\item An \textbf{attention--bone IoU} interpretability metric \emph{with
      controls} (matched vs.\ mismatched vs.\ random), establishing that the
      learned attention anchors on each individual's rigid bone geometry ---
      a control no prior occlusion/attention face-recognition work reports.
\end{enumerate}

% =============================================================================
\section{Related Work}
\label{sec:related}

\textbf{Margin-based identity losses.} ArcFace \cite{deng2019arcface},
CosFace \cite{wang2018cosface}, SphereFace \cite{liu2017sphereface} and
AdaCos \cite{zhang2019adacos} add geometric margins to softmax to enlarge
inter-class separation. They optimise clean identity discrimination only and
provide no mechanism to suppress occlusion or lighting nuisance.

\textbf{Occlusion- and mask-robust recognition.} PDSN \cite{song2019occ}
learns pairwise differential masks to drop corrupted feature elements; FROM
\cite{qiu2021from} predicts feature masks from occlusion patterns; the
masked-face literature \cite{anwar2020masked,wang2021maskinv} retrains on
masked data or inpaints the occluded region. These methods improve occluded
accuracy but (a) typically require an occlusion/mask signal or a learned mask
predictor at inference, and (b) do not \emph{prove}, with a controlled metric,
\emph{where} the surviving identity signal lives. MDIE needs no occlusion
detector at inference (it deploys as a plain ArcFace encoder) and supplies that
proof via attention--bone IoU.

\textbf{Attention/landmark priors for faces.} Part- and attention-based
recognisers \cite{song2019occ} and landmark-supervised attention (LAFS-style)
inject spatial structure, but supervise attention toward generic facial parts
rather than the \emph{rigid bone scaffold} that survives disguise, and report
no per-individual specificity control. RATA supervises attention onto each
face's own detected bone landmarks and validates it against mismatched-face and
random-map nulls.

\textbf{Domain-adversarial training.} The gradient-reversal layer of
\cite{ganin2015dann} underlies our AMD head. To our knowledge, adversarial
disentanglement of \emph{occlusion/lighting type} (rather than dataset domain)
inside a face encoder has not been published.

% =============================================================================
\section{Failure-Mode Benchmark}
\label{sec:bench}

\subsection{Datasets and Modifications}
We build the benchmark on LFW (\texttt{sklearn.fetch\_lfw\_people},
$\geq\!8$ images per identity), with an identity-disjoint $80/20$ split that
yields $173$ training and $44$ held-out (unseen) identities. For each image we
deterministically synthesise nine variants:
\begin{itemize}\itemsep0pt
\item \textbf{clean}: original cropped $112{\times}112$ face.
\item \textbf{disguise\_glasses}: randomised clear/sunglass lenses with rims,
      bridge and temple arms.
\item \textbf{disguise\_mask}: surgical mask over nose and mouth.
\item \textbf{disguise\_cap}: forehead/brow occluding cap.
\item \textbf{occlusion\_random}: random rectangular cut-out, 20--40\% of
      face area.
\item \textbf{low\_light}: $\gamma$-darkening with sensor noise.
\item \textbf{over\_exposure}: blown highlights / bloom.
\item \textbf{harsh\_shadow}: directional sigmoid light/shadow ramp.
\item \textbf{aging}: texture + low-pass jaw smoothing + contrast proxy.
\item \textbf{adversarial}: transferred FGSM, $\varepsilon\!=\!4/255$.
\end{itemize}
The occlusion family is $\{$glasses, mask, cap, random occluder$\}$ and the
lighting family is $\{$low-light, over-exposure, harsh-shadow$\}$.

\subsection{Verification Protocol}
For each (model, modification) cell we sample $3\,000$ positive and $3\,000$
negative pairs from the held-out identities (seed $42$). The reference image is
presented \emph{clean}; the probe carries the modification under test (the same
modification is applied to imposters). We report AUC, EER and TAR at fixed FAR.

\subsection{Comparably-Trained Baselines}
Table~\ref{tab:baselines} reports the four baselines trained in the same regime
as MDIE ($\sim\!13$k LFW images, from scratch) --- the only honest
apples-to-apples comparison. All sit at $0.60$--$0.75$ pooled AUC and collapse
hardest on masks and caps, motivating MDIE.

\begin{table}[t]
\centering
\caption{Comparably-trained baselines on the held-out occlusion/lighting
benchmark (AUC $\uparrow$).}
\label{tab:baselines}
\begin{tabular}{lcccc}
\toprule
Model & clean & occlusion & lighting & pooled \\
\midrule
FaceNet        & 0.664 & 0.591 & 0.571 & 0.603 \\
ArcFace        & 0.703 & 0.629 & 0.609 & 0.641 \\
CosFace        & 0.833 & 0.716 & 0.720 & 0.749 \\
MobileFaceNet  & 0.849 & 0.689 & 0.755 & 0.749 \\
\bottomrule
\end{tabular}
\end{table}

% =============================================================================
\section{The MDIE Encoder}
\label{sec:method}

\subsection{Overview}
MDIE keeps the IR-50 backbone of ArcFace and adds three training components and
one inference-time fusion head:
\begin{enumerate}
\item ArcFace identity head $H_{\text{id}}$ (margin $m\!=\!0.5$, scale
      $s\!=\!64$) on the fused embedding.
\item AMD modification classifier $H_{\text{mod}}$ on a Gradient-Reversal
      Layer (GRL), adversarial weight $\lambda_{\text{amd}}\!=\!0.10$.
\item RATA attention supervised onto rigid-bone landmarks with weight
      $\lambda_{\text{attn}}$.
\item ICCL pair-loss $\mathcal{L}_{\text{iccl}}$ ($\lambda_{\text{iccl}}\!=\!0.50$).
\end{enumerate}
The total objective is
\begin{equation}
\mathcal{L} = \mathcal{L}_{\text{arc}}
            + \lambda_{\text{iccl}}\mathcal{L}_{\text{iccl}}
            + \lambda_{\text{amd}}\mathcal{L}_{\text{amd}}^{\text{GRL}}
            + \lambda_{\text{attn}}\mathcal{L}_{\text{attn}}.
\end{equation}

\subsection{Region-Aware Token Attention (RATA)}
For each face we detect rigid bone landmarks --- brow ridge, orbital rims,
nasal bridge, cheekbones, jaw angles, chin --- the skeletal structure that
survives masks, caps, glasses and lighting. We splat them onto a $14{\times}14$
token grid built from a fused $768$-channel feature map (mid-resolution +
upsampled final stage), giving each landmark \emph{equal mass} so the densely
clustered central T-zone cannot collapse into a single central blob. RATA
injects the prior as an additive attention bias
$\mathrm{softmax}(QK^{\!\top}\!/\sqrt{d} + \lambda M_{\text{region}})V$ and is
supervised by a symmetric (forward+reverse) matching loss
$\mathcal{L}_{\text{attn}}$ with a background-suppression gate. At $7{\times}7$
distinct landmarks collapse into one token; $14{\times}14$ resolves individual
bones, yielding per-face-distinct attention.

\subsection{Inter-Condition Consistency Loss (ICCL)}
For each minibatch we form pairs $(x_i,\tilde{x}_i)$ where
$\tilde{x}_i\!=\!\mathcal{M}_{m_i}(x_i)$. With L2-normalised embeddings, ICCL is
\begin{equation}
\mathcal{L}_{\text{iccl}} = -\frac{1}{B}\sum_{i=1}^{B}
\log\frac{\exp(z_i^{\!\top}\tilde{z}_i/\tau)}
        {\sum_{j\in N(i)} w_{ij}\exp(z_i^{\!\top}\tilde{z}_j/\tau)},
\end{equation}
with $w_{ij}\!=\!2$ when $m_j\!=\!m_i$ (condition-aware hard mining): the
hardest negative is the \emph{same disguise on a different person}.

\subsection{Adversarial Modification Disentanglement (AMD)}
$H_{\text{mod}}$ predicts the modification class from the embedding through a
GRL; the backward pass rescales the gradient by $-\lambda_{\text{amd}}$, so the
encoder is optimised to remove modification information while $H_{\text{mod}}$
recovers it --- a minimax analogous to domain-adversarial training
\cite{ganin2015dann}. Labels are needed only at training time.

\subsection{ArcFace-Compatible Inference}
At test time we drop $H_{\text{id}}$ and $H_{\text{mod}}$. The bone-anchored
attention embedding and the native identity embedding are concatenated and
passed through a learned fusion head ($\mathrm{Linear}(1024\!\to\!512)$ +
BatchNorm), then L2-normalised, producing a \emph{single $512$-d vector in one
forward pass}. Because the ArcFace head and ICCL are trained \emph{through} the
fused embedding, the deployed vector is exactly what is optimised. We verify
(Sec.~\ref{sec:deploy}) that this is a literal ArcFace drop-in: unit-norm,
cosine $=$ dot product, deterministic, no test-time augmentation.

% =============================================================================
\section{Experiments}
\label{sec:results}

\subsection{Setup}
Models train on a single NVIDIA RTX\,3050 (4\,GB) with PyTorch + CUDA mixed
precision; the production bundle re-targets NVIDIA A100-SXM4 via SLURM. All
reported MDIE numbers use the deployed single-$512$-d fused embedding.

\subsection{Main Results and Ablation}
Table~\ref{tab:ablation} reports the family-level AUC. MDIE-full attains
pooled AUC $0.979$ versus $0.749$ for the best comparably-trained baseline
($+0.23$), with occlusion family $0.975$ ($+0.26$) and lighting family $0.980$
($+0.22$). The ablation ladder is \emph{monotone} ---
full $>$ noICCL $\approx$ noAMD $>$ noRATA --- so every component contributes.
AMD and ICCL carry most of the raw AUC; RATA's decisive, measurable
contribution is interpretability (Sec.~\ref{sec:interp}).

\begin{table}[t]
\centering
\caption{MDIE ablation on the held-out benchmark (AUC $\uparrow$, single
$512$-d fused embedding). Best comparably-trained baseline shown for reference.}
\label{tab:ablation}
\begin{tabular}{lcccc}
\toprule
Model & clean & occlusion & lighting & pooled \\
\midrule
\textbf{MDIE-full}  & 0.984 & \textbf{0.975} & \textbf{0.980} & \textbf{0.979} \\
MDIE-noICCL         & 0.981 & 0.970 & 0.975 & 0.974 \\
MDIE-noAMD          & 0.981 & 0.970 & 0.975 & 0.974 \\
MDIE-noRATA         & 0.977 & 0.969 & 0.973 & 0.972 \\
\midrule
Best baseline       & 0.849 & 0.716 & 0.755 & 0.749 \\
\bottomrule
\end{tabular}
\end{table}

Per-modification AUC/EER for MDIE-full: clean $0.984/5.6\%$, glasses
$0.981/6.6\%$, mask $0.966/9.2\%$ (hardest), cap $0.970/7.9\%$, occluder
$0.981/6.3\%$, low-light $0.979/6.7\%$, over-exposure $0.980/7.0\%$,
harsh-shadow $0.981/6.0\%$, aging $0.982/6.4\%$, adversarial $0.983/5.8\%$.
Even the hardest case (mask) clears the best baseline by $+0.28$ AUC.

\subsection{Transfer to Real Worn Occlusions}
Table~\ref{tab:real} evaluates the \emph{same} deployed model on public
protocols. MFR2 (real masks) and MeGlass (real eyeglasses) are the acid test
that the bone-anchored embedding transfers off the synthetic training
distribution. MDIE wins every comparably-trained comparison by $+0.15$ to
$+0.34$ AUC, and remains ahead even on off-niche cross-age sets (CALFW,
AgeDB-30) where the baselines sit near chance.

\begin{table}[t]
\centering
\caption{Real-benchmark transfer (AUC $\uparrow$). Best comparably-trained
baseline vs.\ MDIE (same deployed $512$-d model).}
\label{tab:real}
\begin{tabular}{lccc}
\toprule
Benchmark & pairs & best baseline & \textbf{MDIE} \\
\midrule
MFR2 (real masks)    & 848   & 0.677 & \textbf{0.826} \\
MeGlass (real glasses)& 3000  & 0.706 & \textbf{0.926} \\
CALFW (cross-age)    & 6000  & 0.516 & \textbf{0.599} \\
AgeDB-30             & 6000  & 0.500 & \textbf{0.683} \\
\bottomrule
\end{tabular}
\end{table}

We do not claim to beat production InsightFace \texttt{w600k\_r50} (trained on
$\sim\!17$M faces) on raw accuracy; the honest claim is that MDIE is the most
robust recogniser \emph{in the comparably-trained regime} on the
occlusion/lighting niche, and that the framework re-trains at scale on the
provided A100 bundle.

\subsection{Attention--Bone IoU Interpretability}
\label{sec:interp}
On the $44$ held-out identities we binarise the learned attention and each
face's own rigid-bone target to their top cells and measure IoU, with two
controls (Table~\ref{tab:iou}). The \textbf{matched $\gg$ mismatched $\gg$
random} ordering holds across top-$10/15/20/25\%$ thresholds (matched$-$mismatched
gap $0.36$--$0.41$), Mann--Whitney $p\!=\!1.6\!\times\!10^{-15}$. Every
anatomical group is attended above its uniform share (nose-bridge $13.2\times$,
orbital rim $5.6\times$, cheekbone $4.9\times$, jaw/chin $3.2\times$, brow
$3.0\times$). The attention is thus anchored on \emph{each individual's} bone
geometry --- not a shared face-layout template, not chance.
Fig.~\ref{fig:iou} (\texttt{attention\_bone\_iou.pdf}).

\begin{table}[t]
\centering
\caption{Attention--bone IoU with controls ($44$ held-out identities,
top-$15\%$ threshold).}
\label{tab:iou}
\begin{tabular}{lc}
\toprule
Condition & IoU $\uparrow$ \\
\midrule
matched (face's own bones)        & \textbf{0.694} $\pm$ 0.090 \\
mismatched (a different face's bones) & 0.280 $\pm$ 0.100 \\
random-attention null             & 0.077 \\
\bottomrule
\end{tabular}
\end{table}

\subsection{ArcFace-Compatible Deployment}
\label{sec:deploy}
We verify the deployed embedding is a literal ArcFace drop-in
(\texttt{inference\_compat\_proof.json}): output shape $[B,512]$; L2 unit-norm
(max error $6\!\times\!10^{-8}$); \texttt{encode\_verify} equals
\texttt{encode} (single forward, zero error, no test-time augmentation);
cosine $=$ dot product on the unit sphere (so any ArcFace / FAISS
inner-product index applies unchanged); deterministic in eval mode. A masked
probe still matches its clean gallery image (cosine $0.876$) above an imposter
(cosine $0.824$).

% =============================================================================
\section{Discussion}
\label{sec:discussion}

\textbf{Why the design works.} AMD removes the occlusion/lighting shortcut as a
confounder; ICCL requires (not merely hopes for) clean--modified invariance;
RATA biases attention toward the rigid bone scaffold and \emph{proves} the
surviving signal is anatomical rather than a dataset shortcut. The fusion head
packages all of this as a single cosine-matched vector, so robustness is gained
at \emph{zero} inference overhead and no label dependence.

\textbf{Scale, not architecture, is the ceiling.} We train on $\sim\!13$k
images; production references use $\sim\!17$M. We therefore frame the
contribution as the occlusion/lighting-niche win, the bone-anchored
interpretability control, and the ArcFace-compatible deployment --- all
retrainable at scale on the included PARAM A100 SLURM bundle.

\textbf{Limitations.} (i) Training modifications are protocol-faithful but
synthetic; the real MFR2/MeGlass results address occlusion transfer, but a
\emph{real} low-light/shadow benchmark remains future work. (ii) Reported AUCs
are single-run; multi-seed confidence intervals are left for the scale-up.
(iii) RATA's contribution is interpretability rather than raw AUC; on pooled
AUC the variants lie within $\sim\!0.007$.

% =============================================================================
\section{Conclusion}
\label{sec:concl}
We re-scoped modification-robust face recognition to the demonstrable security
niche of worn occlusion and adverse lighting, and showed that an inexpensive
combination of adversarial disentanglement, consistency contrastive learning,
and bone-anchored attention --- merged into a single ArcFace-compatible
$512$-d embedding --- lifts pooled AUC from the $0.749$ baseline regime to
$0.979$, transfers to real masked and eyeglass benchmarks, and is provably
anchored on each individual's rigid bone structure. The full codebase, figures,
and a one-command reproduction (plus an A100 SLURM bundle) are released.

% =============================================================================
\section*{Reproducibility}
\noindent\texttt{python -m research\_v2.src.run\_stage2 --epochs 12 --ablation} \\
\noindent\texttt{python research\_v2/scripts/attention\_bone\_iou.py} \\
\noindent\texttt{python -m research\_v2.src.eval.run\_real\_benchmarks} \\
\noindent\texttt{python research\_v2/scripts/inference\_compat\_proof.py}

\noindent Seeds and library versions are pinned in
\texttt{research\_v2/src/config.py}; the A100 pipeline is
\texttt{hpc/slurm\_full\_pipeline.sh}.

% =============================================================================
\bibliographystyle{IEEEtran}
\begin{thebibliography}{9}
\bibitem{deng2019arcface}
J.~Deng, J.~Guo, N.~Xue, and S.~Zafeiriou,
``ArcFace: Additive angular margin loss for deep face recognition,'' CVPR, 2019.

\bibitem{wang2018cosface}
H.~Wang \emph{et al.},
``CosFace: Large margin cosine loss for deep face recognition,'' CVPR, 2018.

\bibitem{schroff2015facenet}
F.~Schroff, D.~Kalenichenko, and J.~Philbin,
``FaceNet: A unified embedding for face recognition and clustering,'' CVPR, 2015.

\bibitem{liu2017sphereface}
W.~Liu \emph{et al.},
``SphereFace: Deep hypersphere embedding for face recognition,'' CVPR, 2017.

\bibitem{zhang2019adacos}
X.~Zhang \emph{et al.},
``AdaCos: Adaptively scaling cosine logits for effectively learning deep
face representations,'' CVPR, 2019.

\bibitem{song2019occ}
L.~Song \emph{et al.},
``Occlusion robust face recognition based on mask learning with pairwise
differential siamese network (PDSN),'' ICCV, 2019.

\bibitem{qiu2021from}
H.~Qiu \emph{et al.},
``End2End occluded face recognition by masking corrupted features (FROM),''
IEEE TPAMI, 2021.

\bibitem{anwar2020masked}
A.~Anwar and A.~Raychowdhury,
``Masked face recognition for secure authentication,'' arXiv:2008.11104, 2020.

\bibitem{wang2021maskinv}
Z.~Wang \emph{et al.},
``Mask-invariant face recognition via masked face inpainting,''
IEEE Access, 2021.

\bibitem{ganin2015dann}
Y.~Ganin and V.~Lempitsky,
``Unsupervised domain adaptation by backpropagation,'' ICML, 2015.
\end{thebibliography}

\end{document}
